\documentclass[a4paper,12pt]{ltjsarticle}
\usepackage{luatexja}
\usepackage{amsmath,amssymb,amsthm}
\usepackage{tikz}
\usepackage{tikz-cd}
\usetikzlibrary{positioning,arrows.meta,decorations.pathmorphing,shapes.geometric}
\usepackage{xcolor}
\usepackage[colorlinks=true,linkcolor=blue,citecolor=blue,urlcolor=blue]{hyperref}
\usepackage{tcolorbox}
\tcbuselibrary{theorems,skins,breakable}

% 定理環境
\theoremstyle{definition}
\newtheorem{definition}{定義}[section]
\newtheorem{theorem}[definition]{定理}
\newtheorem{proposition}[definition]{命題}
\newtheorem{lemma}[definition]{補題}
\newtheorem{example}[definition]{例}
\newtheorem{remark}[definition]{注意}

% カラーボックス
\newtcolorbox{keypoint}[1][]{
  colback=blue!5!white,
  colframe=blue!75!black,
  fonttitle=\bfseries,
  title=重要ポイント,
  #1
}

\newtcolorbox{motivation}[1][]{
  colback=green!5!white,
  colframe=green!75!black,
  fonttitle=\bfseries,
  title=数学的動機,
  #1
}

% タイトル情報
\title{σ-代数の構造塔理論とフィルトレーション\\
{\large 確率論への応用を見据えた学部生向け解説}}
\author{
  Lean 4 形式化コード生成：CODEX (OpenAI)\\
  数学的解説とTeX生成：Claude Code (Anthropic)
}
\date{\today}

\begin{document}

\maketitle

\begin{abstract}
本稿は、確率論の基礎となるσ-代数とフィルトレーションの概念を、「構造塔」という統一的な視点から解説する。
構造塔理論は、様々な数学的対象における階層構造を抽象化した枠組みであり、
σ-代数の包含関係や時間発展するフィルトレーションを自然に記述できる。
本稿では、Lean 4による形式化を通じて、これらの概念の厳密な定義と性質を明確にし、
停止時間などの確率論的概念との関連を示す。
\end{abstract}

\tableofcontents

\section{導入：なぜ構造塔なのか}

\subsection{数学における階層構造}

\begin{motivation}
数学の多くの分野で、「より小さな構造」から「より大きな構造」へという階層的な包含関係が現れる：
\begin{itemize}
  \item \textbf{位相空間論}：粗い位相 $\subseteq$ 細かい位相
  \item \textbf{測度論}：小さなσ-代数 $\subseteq$ 大きなσ-代数
  \item \textbf{代数学}：部分群の鎖 $H_1 \subseteq H_2 \subseteq \cdots \subseteq G$
  \item \textbf{確率論}：時間とともに増える情報（フィルトレーション）
\end{itemize}
これらを統一的に扱う枠組みが「構造塔」(Structure Tower) である。
\end{motivation}

\subsection{構造塔の直感的イメージ}

構造塔とは、次のような「階層構造」を持つ数学的対象である：

\begin{center}
\begin{tikzpicture}[
  layer/.style={draw, rectangle, minimum width=8cm, minimum height=0.8cm, fill=blue!20},
  element/.style={circle, fill=red!70, inner sep=2pt}
]
  % 層を描画
  \node[layer] (L4) at (0,4) {第4層（最も大きい）};
  \node[layer] (L3) at (0,3) {第3層};
  \node[layer] (L2) at (0,2) {第2層};
  \node[layer] (L1) at (0,1) {第1層};
  \node[layer] (L0) at (0,0) {第0層（最も小さい）};

  % 要素の例
  \node[element, label=right:$x_1$] at (-3,0.5) {};
  \node[element, label=right:$x_2$] at (1,1.5) {};
  \node[element, label=right:$x_3$] at (-1,2.5) {};

  % 矢印
  \draw[-Latex, thick, blue] (6,0.5) -- (6,3.5) node[midway, right] {包含関係};

  % 注釈
  \node[align=left, text width=7cm] at (0,-1.5) {
    各要素 $x$ は、ある「最小の層」に初めて現れ、\\
    それより上の層にはすべて含まれる
  };
\end{tikzpicture}
\end{center}

\begin{keypoint}
構造塔の3つの本質的性質：
\begin{enumerate}
  \item \textbf{被覆性}：すべての要素は、どこかの層に含まれる
  \item \textbf{単調性}：下の層 $\subseteq$ 上の層（階層構造）
  \item \textbf{最小層の存在}：各要素が「初めて現れる」最小の層が存在する
\end{enumerate}
\end{keypoint}

\section{構造塔の形式的定義}

\begin{definition}[構造塔（最小層付き）]\label{def:structure-tower}
前順序集合 $(I, \leq)$ で添字付けられた集合 $X$ 上の\textbf{構造塔} (Structure Tower with Minimal Layer) とは、
以下のデータと公理からなる：

\textbf{データ：}
\begin{itemize}
  \item 層関数 $\mathrm{layer}: I \to \mathcal{P}(X)$（各添字 $i$ に対し部分集合 $\mathrm{layer}(i) \subseteq X$）
  \item 最小層関数 $\mathrm{minLayer}: X \to I$（各要素の最小層）
\end{itemize}

\textbf{公理：}
\begin{enumerate}
  \item \textbf{被覆性} (covering)：
  \[
    \forall x \in X,\ \exists i \in I,\ x \in \mathrm{layer}(i)
  \]

  \item \textbf{単調性} (monotone)：
  \[
    \forall i, j \in I,\ i \leq j \implies \mathrm{layer}(i) \subseteq \mathrm{layer}(j)
  \]

  \item \textbf{最小層への帰属} (minLayer\_mem)：
  \[
    \forall x \in X,\ x \in \mathrm{layer}(\mathrm{minLayer}(x))
  \]

  \item \textbf{最小性} (minLayer\_minimal)：
  \[
    \forall x \in X,\ \forall i \in I,\ x \in \mathrm{layer}(i) \implies \mathrm{minLayer}(x) \leq i
  \]
\end{enumerate}
\end{definition}

\begin{remark}
Lean 4での形式化では、この定義は次のように記述される：
\begin{verbatim}
structure StructureTowerMin (X : Type*) (I : Type*) [Preorder I] where
  layer : I → Set X
  covering : ∀ x : X, ∃ i : I, x ∈ layer i
  monotone : ∀ {i j : I}, i ≤ j → layer i ⊆ layer j
  minLayer : X → I
  minLayer_mem : ∀ x, x ∈ layer (minLayer x)
  minLayer_minimal : ∀ x i, x ∈ layer i → minLayer x ≤ i
\end{verbatim}
\end{remark}

\section{σ-代数の構造塔}

\subsection{σ-代数の復習}

\begin{definition}[σ-代数]\label{def:sigma-algebra}
集合 $\Omega$ 上の\textbf{σ-代数} (sigma-algebra) $\mathcal{F}$ とは、
$\Omega$ の部分集合族で以下を満たすもの：
\begin{enumerate}
  \item $\Omega \in \mathcal{F}$
  \item $A \in \mathcal{F} \implies A^c \in \mathcal{F}$（補集合で閉じている）
  \item $(A_n)_{n \in \mathbb{N}} \subseteq \mathcal{F} \implies \bigcup_{n=1}^\infty A_n \in \mathcal{F}$（可算和で閉じている）
\end{enumerate}
\end{definition}

\begin{example}[基本的なσ-代数]
\begin{itemize}
  \item \textbf{自明なσ-代数}：$\{\emptyset, \Omega\}$（最も粗い）
  \item \textbf{ベキ集合}：$\mathcal{P}(\Omega) = 2^\Omega$（最も細かい）
  \item \textbf{Borelσ-代数}：$\mathcal{B}(\mathbb{R})$（開集合から生成）
\end{itemize}
\end{example}

\subsection{σ-代数の順序構造}

σ-代数には、包含関係による\textbf{半順序}が自然に入る：

\begin{definition}[σ-代数の順序]
同じ $\Omega$ 上の2つのσ-代数 $\mathcal{F}_1, \mathcal{F}_2$ に対し、
\[
  \mathcal{F}_1 \leq \mathcal{F}_2 \stackrel{\text{def}}{\iff}
  \forall A,\ A \text{ が } \mathcal{F}_1\text{-可測} \implies A \text{ が } \mathcal{F}_2\text{-可測}
\]
と定義する。これは $\mathcal{F}_2$ の方が「細かい」（より多くの集合を測れる）ことを意味する。
\end{definition}

\begin{center}
\begin{tikzpicture}[node distance=2cm]
  \node[draw, ellipse, minimum width=1.5cm, minimum height=1cm] (F1) {$\mathcal{F}_1$};
  \node[draw, ellipse, minimum width=3cm, minimum height=2cm, right=of F1] (F2) {$\mathcal{F}_2$};
  \node[draw, ellipse, minimum width=5cm, minimum height=3.5cm, right=of F2] (F3) {$\mathcal{F}_3$};

  \draw[-Latex, thick] (F1) -- (F2) node[midway, above] {$\leq$};
  \draw[-Latex, thick] (F2) -- (F3) node[midway, above] {$\leq$};

  \node[below=0.5cm of F1] {粗い};
  \node[below=0.5cm of F3] {細かい};
\end{tikzpicture}
\end{center}

\subsection{σ-代数の構造塔の構成}

\begin{theorem}[σ-代数構造塔の存在]\label{thm:sigma-algebra-tower}
任意の集合 $\Omega$ に対し、次のような構造塔が存在する：
\begin{itemize}
  \item 要素空間 $X = \mathcal{P}(\Omega)$（$\Omega$ の部分集合全体）
  \item 添字集合 $I = \mathrm{MeasurableSpace}(\Omega)$（$\Omega$ 上のσ-代数全体）
  \item 層関数：$\mathrm{layer}(\mathcal{F}) = \{ A \subseteq \Omega \mid A \text{ は } \mathcal{F}\text{-可測}\}$
  \item 最小層：$\mathrm{minLayer}(A) = $ 「$A$ を可測にする最小のσ-代数」$= \sigma(\{A\})$
\end{itemize}
\end{theorem}

\begin{proof}[証明のスケッチ]
各公理を検証する：

\textbf{(被覆性)}：任意の集合 $A \subseteq \Omega$ に対し、ベキ集合σ-代数 $\mathcal{F} = \mathcal{P}(\Omega)$ では
すべての集合が可測なので、$A \in \mathrm{layer}(\mathcal{P}(\Omega))$。

\textbf{(単調性)}：$\mathcal{F}_1 \leq \mathcal{F}_2$ ならば、
$A$ が $\mathcal{F}_1$-可測ならば $\mathcal{F}_2$-可測なので、
$\mathrm{layer}(\mathcal{F}_1) \subseteq \mathrm{layer}(\mathcal{F}_2)$。

\textbf{(最小層への帰属)}：生成σ-代数 $\sigma(\{A\})$ の定義により、$A$ は $\sigma(\{A\})$-可測。

\textbf{(最小性)}：$A$ が $\mathcal{F}$-可測ならば、$\sigma(\{A\}) \leq \mathcal{F}$（生成の最小性）。
\end{proof}

\begin{center}
\begin{tikzpicture}[
  sigma/.style={draw, rectangle, rounded corners, minimum width=7cm, minimum height=0.7cm, fill=green!20},
  set/.style={draw, circle, fill=red!60, inner sep=1pt}
]
  % σ-代数の階層
  \node[sigma] (S4) at (0,4) {$\mathcal{P}(\Omega)$：すべての集合が可測};
  \node[sigma] (S3) at (0,3) {$\mathcal{F}_3$：より細かいσ-代数};
  \node[sigma] (S2) at (0,2) {$\mathcal{F}_2$};
  \node[sigma] (S1) at (0,1) {$\mathcal{F}_1$：粗いσ-代数};
  \node[sigma] (S0) at (0,0) {$\{\emptyset, \Omega\}$：最も粗い};

  % 集合の例
  \node[set, label=left:$A$] at (-3,1.5) {};
  \draw[-Latex, red, thick] (-2.8,1.5) -- (-0.5,1.5) node[midway, above, text=black] {初めて可測};

  % 注釈
  \node[align=left, text width=7cm] at (0,-1.5) {
    集合 $A$ が初めて可測になるσ-代数が\\
    $\mathrm{minLayer}(A) = \sigma(\{A\})$（生成σ-代数）
  };
\end{tikzpicture}
\end{center}

\section{フィルトレーション：時間発展する情報}

\subsection{確率論におけるフィルトレーション}

\begin{motivation}
確率論では、時間とともに「得られる情報」が増えていく状況を扱う：
\begin{itemize}
  \item コイン投げ：1回目、2回目、…と結果が明らかになる
  \item 株価の観測：時刻 $t$ までの価格履歴
  \item ブラウン運動：時刻 $t$ までの軌道の情報
\end{itemize}
この「情報の増加」を数学的に表現するのが\textbf{フィルトレーション}である。
\end{motivation}

\begin{definition}[フィルトレーション]\label{def:filtration}
確率空間 $(\Omega, \mathcal{F}, P)$ 上の\textbf{フィルトレーション} (filtration) とは、
σ-代数の単調増加列
\[
  \mathcal{F}_0 \subseteq \mathcal{F}_1 \subseteq \mathcal{F}_2 \subseteq \cdots \subseteq \mathcal{F}
\]
である。$\mathcal{F}_n$ は「時刻 $n$ までに得られる情報」を表すσ-代数と解釈される。
\end{definition}

\begin{center}
\begin{tikzpicture}[
  time/.style={draw, circle, minimum size=1cm, fill=blue!30},
  sigma/.style={draw, rectangle, minimum width=2cm, minimum height=0.8cm, fill=yellow!30}
]
  % 時間軸
  \draw[-Latex, ultra thick] (-0.5,0) -- (10,0) node[right] {時間 $t$};

  % 各時刻
  \foreach \t/\x in {0/0, 1/2, 2/4, 3/6, n/9} {
    \node[time] at (\x,0) {$t=\t$};
    \node[sigma] (F\t) at (\x,1.5) {$\mathcal{F}_\t$};
    \draw[-Latex] (\x,0.5) -- (F\t);
  }

  % 包含関係
  \draw[-Latex, thick, red] (F0) -- (F1) node[midway, above] {$\subseteq$};
  \draw[-Latex, thick, red] (F1) -- (F2) node[midway, above] {$\subseteq$};
  \draw[-Latex, thick, red] (F2) -- (F3) node[midway, above] {$\subseteq$};
  \draw[dashed, thick] (F3) -- (Fn);

  \node[align=center, text width=8cm] at (5,-2) {
    時間が進むにつれ、より多くの事象が「観測可能」になる\\
    $\mathcal{F}_n$ = 時刻 $n$ までに決定できる事象全体
  };
\end{tikzpicture}
\end{center}

\subsection{フィルトレーションから構造塔へ}

\begin{proposition}[フィルトレーションは構造塔]\label{prop:filtration-to-tower}
フィルトレーション $(\mathcal{F}_n)_{n \in \mathbb{N}}$ が次の条件を満たすとき：
\[
  \forall A \subseteq \Omega,\ \exists n \in \mathbb{N},\ A \text{ は } \mathcal{F}_n\text{-可測}
  \quad \text{(最終的被覆条件)}
\]
このフィルトレーションから構造塔が構成できる：
\begin{itemize}
  \item 要素空間：$X = \mathcal{P}(\Omega)$
  \item 添字集合：$I = \mathbb{N}$（時刻）
  \item 層関数：$\mathrm{layer}(n) = \{ A \mid A \text{ は } \mathcal{F}_n\text{-可測}\}$
  \item 最小層：$\mathrm{minLayer}(A) = \min\{ n \mid A \in \mathcal{F}_n \}$（初めて可測になる時刻）
\end{itemize}
\end{proposition}

\begin{proof}
\textbf{(被覆性)}：最終的被覆条件より従う。

\textbf{(単調性)}：フィルトレーションの定義 $\mathcal{F}_m \subseteq \mathcal{F}_n$ ($m \leq n$) より、
$\mathrm{layer}(m) \subseteq \mathrm{layer}(n)$。

\textbf{(最小層)}：$\min$ の存在は well-ordering principle より従う。
最小性は $\min$ の定義より明らか。
\end{proof}

\begin{example}[コイン投げのフィルトレーション]
コインを無限回投げる：$\Omega = \{H, T\}^\mathbb{N}$。
\begin{itemize}
  \item $\mathcal{F}_0 = \{\emptyset, \Omega\}$：何も知らない
  \item $\mathcal{F}_1$：1回目の結果で決まる事象全体
  \item $\mathcal{F}_2$：1,2回目の結果で決まる事象全体
  \item $\vdots$
  \item $\mathcal{F}_n$：最初の $n$ 回で決まる事象全体
\end{itemize}

集合 $A = $ 「1回目が表」の $\mathrm{minLayer}(A) = 1$（1回目で初めて決定できる）。
\end{example}

\section{停止時間と最小層の関係}

\subsection{停止時間の定義}

\begin{definition}[停止時間]\label{def:stopping-time}
フィルトレーション $(\mathcal{F}_n)$ に対し、確率変数 $\tau: \Omega \to \mathbb{N} \cup \{\infty\}$ が
\textbf{停止時間} (stopping time) であるとは：
\[
  \forall n \in \mathbb{N},\ \{\omega \in \Omega \mid \tau(\omega) \leq n\} \in \mathcal{F}_n
\]
つまり、「時刻 $n$ に停止したかどうか」が時刻 $n$ の情報で決定できる。
\end{definition}

\begin{example}[基本的な停止時間]
\begin{itemize}
  \item \textbf{定時刻}：$\tau(\omega) = k$（常に時刻 $k$）
  \item \textbf{初到達時刻}：$\tau = \inf\{ n \geq 0 \mid X_n \in B \}$（集合 $B$ に初めて入る時刻）
  \item \textbf{初通過時刻}：$\tau = \inf\{ n \geq 0 \mid X_n \geq a \}$（閾値 $a$ を初めて超える時刻）
\end{itemize}
\end{example}

\subsection{停止時間と構造塔}

\begin{center}
\begin{tikzpicture}[
  state/.style={draw, circle, minimum size=0.6cm, fill=orange!40},
  event/.style={draw, ellipse, minimum width=2cm, minimum height=0.8cm, fill=purple!20}
]
  % 状態空間（簡略化）
  \node[state] (w1) at (0,3) {$\omega_1$};
  \node[state] (w2) at (2,3) {$\omega_2$};
  \node[state] (w3) at (4,3) {$\omega_3$};
  \node[state] (w4) at (6,3) {$\omega_4$};

  % 停止時刻
  \draw[-Latex, red, ultra thick] (w1) -- (0,1) node[below] {$\tau(\omega_1)=2$};
  \draw[-Latex, red, ultra thick] (w2) -- (2,0.5) node[below] {$\tau(\omega_2)=1$};
  \draw[-Latex, red, ultra thick] (w3) -- (4,0) node[below] {$\tau(\omega_3)=0$};
  \draw[-Latex, red, ultra thick] (w4) -- (6,1.5) node[below] {$\tau(\omega_4)=3$};

  % 事象
  \node[event] at (-1,0) {$A_0$};
  \node[event] at (1,0.5) {$A_1$};
  \node[event] at (3,1) {$A_2$};
  \node[event] at (5,1.5) {$A_3$};

  % 説明
  \node[align=left, text width=8cm] at (3,-1.5) {
    各 $\omega$ に対し、「$\omega$ に関する決定的事象が\\
    初めて観測できる時刻」が $\tau(\omega)$\\
    $\iff$ その事象の $\mathrm{minLayer}$（構造塔的解釈）
  };
\end{tikzpicture}
\end{center}

\begin{keypoint}
\textbf{停止時間と最小層の関係：}
停止時間 $\tau(\omega)$ は、次のように解釈できる：
\[
  \tau(\omega) = \mathrm{minLayer}_\text{filtration}(\text{「$\omega$ を識別する事象」})
\]
つまり、$\omega$ に関する情報が「初めて可測になる時刻」が $\tau(\omega)$ である。

これは構造塔の $\mathrm{minLayer}$ の確率論的解釈に他ならない。
\end{keypoint}

\subsection{形式化における課題}

\begin{remark}[Lean 4 での表現]
現在の形式化では：
\begin{verbatim}
theorem stopping_time_as_minLayer (ℱ : SigmaAlgebraFiltration (Ω := Ω))
    (τ : StoppingTime ℱ) (ω : Ω) :
    -- τ(ω) は ω に関する何らかの事象の minLayer
    True := by
  trivial  -- TODO: 厳密な定式化と証明
\end{verbatim}

この関係を厳密に定式化するには、以下が必要：
\begin{itemize}
  \item $\omega$ に対応する「決定的事象」の構成
  \item その事象の可測性と最小層の計算
  \item 停止時間の定義との整合性の証明
\end{itemize}
これは今後の研究課題である。
\end{remark}

\section{構造塔の他の応用例}

σ-代数以外にも、構造塔は様々な数学的対象に適用できる：

\begin{example}[位相空間の構造塔]
位相空間 $(X, \tau)$ において、開集合族の包含関係：
\begin{itemize}
  \item 粗い位相 $\tau_1 \subseteq \tau_2$ 細かい位相
  \item $\mathrm{layer}(\tau) = $ 「位相 $\tau$ での開集合全体」
  \item $\mathrm{minLayer}(U) = $ 「$U$ を開集合にする最小の位相」$= \{ \emptyset, U, U^c, X \}$
\end{itemize}
\end{example}

\begin{example}[部分群の構造塔]
群 $G$ の部分群の包含関係：
\begin{itemize}
  \item $H_1 \subseteq H_2 \subseteq \cdots \subseteq G$
  \item $\mathrm{layer}(H) = H$（部分群そのもの）
  \item $\mathrm{minLayer}(g) = \langle g \rangle$（$g$ で生成される巡回部分群）
\end{itemize}
\end{example}

\begin{example}[主イデアルの構造塔]
整数環 $\mathbb{Z}$ の主イデアル：
\begin{itemize}
  \item $(12) \subseteq (6) \subseteq (3) \subseteq (1) = \mathbb{Z}$
  \item 約数関係による包含：$a \mid b \iff (b) \subseteq (a)$
  \item $\mathrm{minLayer}(n) = (n)$（$n$ 自身のイデアル）
\end{itemize}
\end{example}

\begin{center}
\begin{tikzpicture}[
  box/.style={draw, rounded corners, minimum width=3cm, minimum height=1.5cm, fill=cyan!20, align=center}
]
  \node[box] (sigma) at (0,0) {σ-代数\\フィルトレーション};
  \node[box] (top) at (5,0) {位相空間\\開集合の階層};
  \node[box] (alg) at (2.5,-3) {代数構造\\部分群・イデアル};

  \node[draw, ellipse, minimum width=4cm, minimum height=2cm, dashed, ultra thick, red] at (2.5,-1) {};
  \node[text=red, font=\Large\bfseries] at (2.5,-4.5) {構造塔理論で統一};

  \draw[-Latex, thick, blue] (sigma) -- (2.5,-1);
  \draw[-Latex, thick, blue] (top) -- (2.5,-1);
  \draw[-Latex, thick, blue] (alg) -- (2.5,-1);
\end{tikzpicture}
\end{center}

\section{まとめと今後の展望}

\subsection{本稿のまとめ}

本稿では、構造塔理論を通じてσ-代数とフィルトレーションを統一的に理解した：

\begin{enumerate}
  \item \textbf{構造塔の定義}：階層構造を持つ数学的対象の抽象化
  \item \textbf{σ-代数の構造塔}：可測性による階層構造の形式化
  \item \textbf{フィルトレーション}：時間発展する情報の構造塔的記述
  \item \textbf{停止時間}：最小層の確率論的解釈
\end{enumerate}

\begin{keypoint}
\textbf{構造塔理論の意義：}
構造塔理論により、以下が可能になる：
\begin{itemize}
  \item 異なる数学的分野の階層構造を統一的に扱う
  \item 共通の性質（被覆性、単調性、最小層）を一度証明すれば、多くの分野で適用可能
  \item Lean 4 による形式化で、抽象理論の正しさを機械的に検証できる
\end{itemize}
\end{keypoint}

\subsection{今後の発展方向}

\begin{enumerate}
  \item \textbf{測度論への拡張}
  \begin{itemize}
    \item 測度空間の構造塔
    \item 可測関数の階層構造
    \item 積分の構造塔的理解
  \end{itemize}

  \item \textbf{マルチンゲール理論}
  \begin{itemize}
    \item 適合過程 (adapted process) の形式化
    \item 条件付き期待値の階層構造
    \item オプショナル停止定理の構造塔的証明
  \end{itemize}

  \item \textbf{確率過程論への応用}
  \begin{itemize}
    \item ブラウン運動のフィルトレーション
    \item ドゥーブの不等式・定理
    \item 確率収束の階層的理解
  \end{itemize}

  \item \textbf{他分野への展開}
  \begin{itemize}
    \item 代数幾何学：層の構造塔
    \item 関数解析：作用素環の階層
    \item 圏論的な一般化
  \end{itemize}
\end{enumerate}

\subsection{学習のための推奨事項}

本稿の内容をより深く理解するために：

\begin{itemize}
  \item \textbf{測度論の基礎}：Halmos「測度論」、伊藤清三「ルベーグ積分入門」
  \item \textbf{確率論}：Durrett「Probability: Theory and Examples」、Williams「Probability with Martingales」
  \item \textbf{Lean 4 形式化}：「Theorem Proving in Lean 4」、Mathlib のドキュメント
  \item \textbf{圏論的視点}：Awodey「Category Theory」、Mac Lane「Categories for the Working Mathematician」
\end{itemize}

\section*{謝辞}

本稿の Lean 4 形式化コードは CODEX (OpenAI) により生成され、
数学的解説および TeX ドキュメントは Claude Code (Anthropic) により作成された。
構造塔理論の着想は Bourbaki の「母構造」(structures mères) の概念に基づいている。

\appendix

\section{Lean 4 コードの全体}

参考のため、\texttt{SigmaAlgebraTower\_Skeleton.lean} の主要部分を掲載する：

\begin{verbatim}
/-- σ-代数を添字とする構造塔 -/
def SigmaAlgebraTower : StructureTowerMin (Set Ω) (MeasurableSpace Ω) where
  layer 𝓕 := {A : Set Ω | @MeasurableSet Ω 𝓕 A}

  covering := by
    intro A
    use ⊤
    trivial

  monotone := by
    intro 𝓕₁ 𝓕₂ h�� A hA
    exact h𝓕 A hA

  minLayer := fun A =>
    MeasurableSpace.generateFrom {A}

  minLayer_mem := by
    intro A
    exact MeasurableSpace.measurableSet_generateFrom (by simp : A ∈ ({A} : Set (Set Ω)))

  minLayer_minimal := by
    intro A 𝓕 hA
    apply MeasurableSpace.generateFrom_le
    intro B hB
    rcases hB with rfl
    exact hA

/-- フィルトレーションから構造塔を構成 -/
noncomputable def FiltrationToTower (ℱ : SigmaAlgebraFiltration (Ω := Ω)) :
    StructureTowerMin (Set Ω) ℕ where
  layer n := {A : Set Ω | @MeasurableSet Ω (ℱ.𝓕 n) A}

  covering := by
    intro A
    obtain ⟨n, hA⟩ := ℱ.covers A
    exact ⟨n, hA⟩

  monotone := by
    intro m n hmn A hA
    exact ℱ.mono m n hmn A hA

  minLayer := fun A => Nat.find (ℱ.covers A)

  minLayer_mem := by
    intro A
    exact Nat.find_spec (ℱ.covers A)

  minLayer_minimal := by
    intro A n hA
    exact Nat.find_min' (ℱ.covers A) hA
\end{verbatim}

\section{記号表}

\begin{center}
\begin{tabular}{|c|l|}
\hline
記号 & 意味 \\
\hline\hline
$\Omega$ & 標本空間（確率論）または台集合 \\
$\mathcal{F}$ & σ-代数 \\
$\mathcal{F}_n$ & 時刻 $n$ でのσ-代数（フィルトレーション） \\
$\mathrm{layer}(i)$ & 添字 $i$ での構造塔の層 \\
$\mathrm{minLayer}(x)$ & 要素 $x$ の最小層 \\
$\sigma(S)$ & 集合族 $S$ から生成されるσ-代数 \\
$\tau$ & 停止時間 \\
$\mathbb{N}$ & 自然数全体 \\
$\mathcal{P}(X)$ & 集合 $X$ のベキ集合 \\
\hline
\end{tabular}
\end{center}

\end{document}
